% !TeX program = pdflatex
% Counterexamples to two conjectures of Graffiti on graph eigenvalues
% Build: pdflatex note.tex (run twice). See BUILD.md.
\documentclass[11pt]{article}

\usepackage[T1]{fontenc}
\usepackage{lmodern}
\usepackage[margin=1.1in]{geometry}
\usepackage{amsmath,amssymb,amsthm}
\usepackage{booktabs}
\usepackage{microtype}
\usepackage[hidelinks]{hyperref}
\usepackage{xurl}

\newtheorem{theorem}{Theorem}[section]
\newtheorem{proposition}[theorem]{Proposition}
\newtheorem{lemma}[theorem]{Lemma}
\theoremstyle{remark}
\newtheorem{remark}[theorem]{Remark}

\newcommand{\Varp}{\operatorname{Var}^{+}}
\newcommand{\dev}{\operatorname{dev}}
\newcommand{\la}{\ell_{\mathrm{a}}}
\newcommand{\ld}{\ell_{\mathrm{d}}}
\newcommand{\D}{\mathrm{D}}
\newcommand{\Lol}{\mathrm{L}}

\title{Counterexamples to two conjectures of Graffiti\\ on graph eigenvalues}
\author{John Erlbacher\thanks{Independent Researcher. Email: \texttt{erlbacher.research@gmail.com}. ORCID: 0009-0003-6851-4139.}}
\date{June 2026}

\begin{document}
\maketitle

\begin{abstract}
Graffiti, Fajtlowicz's conjecture-making program, produced hundreds of
graph-theoretic conjectures, catalogued in \emph{Written on the Wall}. We
refute two of its spectral conjectures, both of which survived the 1990--91
computational attack of Brewster, Dinneen and Faber (exhaustive only through
ten-vertex graphs) and an eight-algorithm search by Roucairol and
Cazenave (2025). Conjecture~143 asserts that in every
connected graph the variance of the positive adjacency eigenvalues is at most
size/(average distance). It fails on dumbbell graphs (two cliques joined by a
path): the smallest counterexample certified here has $39$ vertices under
both conventions for average distance ($37$ under the distinct-pairs
convention alone), and the ratio of the two
sides tends to $2$ along a dumbbell sequence. Conjecture~154 asserts, under
the standard-deviation reading of ``deviation of eigenvalues'', that the
deviation is at most $n$/(average distance). It fails on lollipop graphs, at
$118$ vertices for distinct pairs and $120$ for both conventions, with
unbounded ratio. All counterexamples carry exact rational certificates,
verified in exact arithmetic by two algorithmically independent routes per
result (for Conjecture~154, by two independently written checkers).
\end{abstract}

\section{Introduction}\label{sec:intro}

Graffiti, written by Siemion Fajtlowicz in the mid-1980s, generates
graph-theoretic conjectures on a large scale
\cite{Fajtlowicz1988,DeLaVina2005}. Its conjectures --- inequalities between
graph invariants --- were collected, numbered, and annotated by Fajtlowicz in
a working document known as \emph{Written on the Wall}
\cite{WOW}\footnote{Numbering throughout this note follows \emph{Written on
the Wall} (WoW). The later, separate list Graffiti.pc of DeLaVi\~na has its
own numbering; its conjecture~143 is unrelated to the one treated here.};
the numbering exceeds 800 in the compilation used here, the file archived by
Roucairol and Cazenave as \texttt{wow-july2004.pdf} (its filename suggests
July 2004, but the PDF's own metadata records a creation date of 26 July
2006; we therefore pin the file by its checksum in
Section~\ref{sec:verification} rather than by date). The list
attracted substantial attention. Favaron, Mah\'eo and Sacl\'e
\cite{FMS1993} resolved many of the early spectral conjectures, proving some
and constructing infinite families of counterexamples to others; Brewster,
Dinneen and Faber \cite{BDF1995} ran a computational attack on roughly 200 of
the conjectures and refuted more than 40. The systematic formalization of
such invariant inequalities was carried further in the AutoGraphiX line of
work of Aouchiche, Hansen and coauthors; see the survey \cite{AouchicheHansen2010}.
Most recently, Roucairol and Cazenave applied Monte-Carlo search to spectral
graph conjectures, among them Graffiti's Conjecture~137
\cite{RoucairolCazenave2022}, and then attacked the Graffiti list with a
portfolio of eight search algorithms
\cite{RoucairolCazenave2025}; the results table of
\cite{RoucairolCazenave2025} lists WoW Conjecture~143 as open after searching
graphs and trees up to size~100, and Conjecture~154 as open after searching up
to size~50.

This note refutes Conjecture~143, and refutes Conjecture~154 under the
standard-deviation reading of ``deviation'' fixed in
Section~\ref{sec:154}.

\begin{quote}
\emph{``143. variance of positive eigenvalues $\le$ size / average
distance.''}\\[2pt]
\emph{``154. deviation of eigenvalues $\le$ n / average distance.''}
\end{quote}

\noindent
Here, by the conventions fixed in the WoW preamble (``$G$ always denotes a
graph, $n$ the number of its vertices and eigenvalues (unless it is
explicitly stated) denote eigenvalues of the adjacency matrix of $G$''),
eigenvalues are those of the adjacency matrix; ``size'' is the number $m$ of
edges; and both statements implicitly concern connected graphs, since the
average distance must be defined. Both conjectures predate the 1990--91
computational attack of \cite{BDF1995}: the WoW compilation records them
among the conjectures that survived it. (That attack was exhaustive over
the graphs on at most ten vertices, so survival there carries limited
information: the counterexamples exhibited below have at least $37$
vertices.) Conjecture~143 fails on
\emph{dumbbells} (two cliques joined by a path), with the smallest
counterexample certified here on $n=39$ vertices --- and on $n=37$ vertices under the
distinct-pairs convention for average distance (Section~\ref{sec:143}).
Conjecture~154, under the standard-deviation reading of ``deviation''
discussed in Section~\ref{sec:154}, fails on \emph{lollipops} (a clique with
a pendant path), first within that family at $n=118$ --- at $n=120$ under
both conventions --- and along an infinite sequence; the violation there
reduces to a pure integer inequality, $8mW^2 > n^5(n-1)^2$, in the Wiener
index $W$. In both
cases the failure is structural rather than marginal: along suitable
parameter sequences the ratio of left- to right-hand side tends to $2$
(Proposition~\ref{prop:143asymp}) and to infinity
(Proposition~\ref{prop:154asymp}), respectively. Both refutations are
robust to the standard ambiguity in the term ``average distance'': they hold
under the two readings in use in this literature (see below). A literature
check (most recently June 2026) found no prior or concurrent resolution of
either conjecture; arXiv:2409.18626 \cite{RoucairolCazenave2025} remains at
version~1, listing both as open.

\paragraph{Conventions.}
Let $G$ be a connected graph with $n$ vertices, $m$ edges, adjacency
eigenvalues $\lambda_1 \ge \dots \ge \lambda_n$, and Wiener index
$W = \sum_{\{u,v\}} d(u,v)$ (sum over unordered pairs of distinct vertices).
Two readings of \emph{average distance} occur in the relevant literature:
\[
\ld = \frac{2W}{n(n-1)}
\qquad\text{and}\qquad
\la = \frac{2W}{n^2},
\]
the mean over distinct (ordered or unordered) pairs, and the mean over all
$n^2$ entries of the distance matrix, zero diagonal included. The
distinct-pairs convention $\ld$ is the one used by Favaron, Mah\'eo and
Sacl\'e \cite{FMS1993} and matches WoW's own usage (its comment on
conjecture~536, asserting that Paley graphs have average distance $3/2$,
holds only under $\ld$); the all-entries convention $\la$ is the one
hard-coded in the public search code of \cite{RoucairolCazenave2025}. Since
$\la < \ld$, right-hand sides computed with $\la$ are \emph{larger}, so the
$\la$ reading is the harder one to violate. We refute both readings of both
conjectures; only the smallest certified instances depend on the reading
(Conjecture~143: $n=39$ for both readings, $n=37$ under $\ld$ alone;
Conjecture~154: $n=118$ under $\ld$, $n=120$ for both).
``Variance'' and ``deviation'' are the population variance and population
standard deviation (divide by the number of data points, not by one less);
WoW computes statistical invariants on uniform sample spaces, and the only
deviation routine published in the search-code repository of
\cite{RoucairolCazenave2025} uses the same convention (see
Section~\ref{sec:154}). For
Conjecture~154 the reading of ``deviation'' carries genuine interpretive
risk, which we discuss --- and quantify --- in Section~\ref{subsec:mad}.

\paragraph{Methods and AI-use disclosure.}
The counterexamples were found by Demonstrandum, a verification-first
multi-agent AI search-and-reasoning pipeline built on Claude language
models (Anthropic): agents fixed the
conjecture statements and conventions from the frozen source of
Section~\ref{sec:verification}, proposed candidate graph families, and ran
parameter scans. Certificates were then verified by exact-arithmetic
checkers written independently of the discovery code: for Conjecture~143, a
single clean-room checker that accepts only if two algorithmically
independent exact routes agree; for Conjecture~154, two independently
written checkers, the second authored by an adversarial referee agent
running on a different model family (Codex, GPT-5.5, OpenAI). All
accept paths use integer or rational arithmetic only; floating point appears
solely in advisory cross-checks. Both verification bundles passed mutation
testing (deliberately corrupted certificates are rejected) and an adversarial
automated audit. The text of this note was likewise drafted with the
Claude-based pipeline; the author directed the work and takes responsibility
for its content. Every number appearing in this note was recomputed during
its preparation (June 2026): the certificate checkers and mutation tests
were re-run, the family scans were re-executed, and the headline quantities
were additionally recomputed by separate scripts
(Section~\ref{sec:verification}). The mathematical claims rest on the
certificates and checkers
described in Section~\ref{sec:verification} and on the self-contained proofs
of Propositions~\ref{prop:143asymp} and~\ref{prop:154asymp}.

\section{Conjecture 143: variance of the positive eigenvalues}\label{sec:143}

For a connected graph $G$ let $k = k(G)$ be the number of strictly positive
adjacency eigenvalues, counted with multiplicity, let $\mu$ be their mean,
and let
\[
\Varp(G) \;=\; \frac1k \sum_{i\colon \lambda_i > 0} (\lambda_i - \mu)^2
\]
be their population variance. Conjecture~143 asserts that
$\Varp(G) \le m/\ell$ for every connected $G$, with $\ell$ the average
distance.

For integers $t_1, t_2 \ge 1$ and $p \ge 0$, the \emph{dumbbell}
$\D(t_1,p,t_2)$ is obtained from disjoint cliques $K_{t_1}$ and $K_{t_2}$ and
a path on $p$ internal vertices by joining one vertex of $K_{t_1}$ to one end
of the path and the other end of the path to one vertex of $K_{t_2}$ (for
$p=0$, a single bridge edge). It is connected, with
$n = t_1 + t_2 + p$ and $m = \binom{t_1}{2} + \binom{t_2}{2} + p + 1$.

\begin{theorem}\label{thm:143}
Conjecture~143 is false, under either reading of average distance. The five
dumbbells of Table~\ref{tab:143} satisfy the indicated strict violations; in
particular, on $n=39$ vertices,
\begin{align*}
\Varp\bigl(\D(7,12,20)\bigr) \;=\; 34.31513046745\ldots
\;&>\; \frac{m\,n^2}{2W} = \frac{10647}{311} = 34.2347266881\ldots\\
\;&>\; \frac{m\,n(n-1)}{2W} = \frac{10374}{311},
\end{align*}
so both readings fail at $n = 39$; and on $n=37$ vertices
\[
\Varp\bigl(\D(6,12,19)\bigr) \;=\; 30.31239262097\ldots
\;>\; \frac{m\,n(n-1)}{2W} = \frac{14726}{487} = 30.23819301848\ldots,
\]
so the distinct-pairs reading fails at $n = 37$.
\end{theorem}

\begin{proof}
A finite exact computation, certified as described in
Section~\ref{sec:verification}. The integers $n$, $m$, $W$, $k$ and the
rational right-hand sides $m/\la = mn^2/(2W)$ and $m/\ld = mn(n-1)/(2W)$ are
elementary exact counts. For $\Varp$, the certificate provides a rational
enclosure of width below $10^{-19}$, obtained by exact real-root isolation of
the integer characteristic polynomial (no floating point on the accept
path); the margin columns of Table~\ref{tab:143} are decimal truncations of
certified exact rational lower bounds on $\Varp - m/\ell$. The clean-room
checker accepts the certificate only after two algorithmically independent
exact routes certify every claimed inequality
(Section~\ref{sec:verification}).
\end{proof}

\begin{table}[ht]
\centering\small
\setlength{\tabcolsep}{3pt}
\begin{tabular}{@{}lrrrrlllll@{}}
\toprule
$G$ & $n$ & $m$ & $W$ & $k$ & $\Varp(G)$ &
$m/\la$ & margin & $m/\ld$ & margin \\
\midrule
$\D(20,8,20)$ & 48 & 389 & 6568 & 6 & 69.6349632447 & $56016/821$  & $+1.405974206$ & $54849/821$  & $+2.827411478$ \\
$\D(8,12,20)$ & 40 & 231 & 5372 & 8 & 34.8927644618 & $46200/1343$ & $+0.492168780$ & $45045/1343$ & $+1.352183672$ \\
$\D(7,12,20)$ & 39 & 224 & 4976 & 8 & 34.3151304675 & $10647/311$  & $+0.080403779$ & $10374/311$  & $+0.958217284$ \\
$\D(6,12,20)$ & 38 & 218 & 4581 & 8 & 33.9573790158 & $157396/4581$ & ---           & $153254/4581$ & $+0.503111388$ \\
$\D(6,12,19)$ & 37 & 199 & 4383 & 8 & 30.3123926210 & $272431/8766$ & ---           & $14726/487$  & $+0.074199602$ \\
\bottomrule
\end{tabular}
\caption{Certified dumbbell counterexamples to Conjecture~143. $\Varp(G)$ is
the certified value rounded to ten decimals (the certificate encloses it in a
rational interval of width $<10^{-19}$). The right-hand sides
$m/\la = mn^2/(2W)$ and $m/\ld = mn(n-1)/(2W)$ are exact. Each margin is a
decimal truncation of a certified exact rational lower bound on
$\Varp(G) - m/\ell$; a dash means the corresponding reading is \emph{not}
violated by that instance --- certified as well: the checker verifies that
the certified enclosure of $\Varp(G)$ lies strictly below that right-hand
side.}
\label{tab:143}
\end{table}

\begin{remark}[Minimality within the family; smaller graphs]\label{rem:143min}
A floating-point scan of the entire dumbbell family with $n \le 48$ (all
$t_1 \ge 1$, $t_2 \ge t_1$, $p \ge 0$) indicates that $\D(7,12,20)$ is the
family-minimal violator for the $\la$ reading and $\D(6,12,19)$ for the
$\ld$ reading, with no dumbbell on $n \le 36$ vertices coming within $0.3$
of either bound (the closest approach, by $\D(6,12,18)$ on $n = 36$
vertices, falls short by about $0.355$) --- far above floating-point error,
though this scan, unlike
Table~\ref{tab:143}, is not an exact certificate. No minimality over
\emph{all} connected graphs is claimed: a simulated-annealing probe at
$n = 35$--$38$ found no smaller violator, but that evidence is heuristic.
The prior search of \cite{RoucairolCazenave2025} covered sizes up to~100;
the counterexamples were found not beyond the searched range but inside it,
in a region apparently unfavourable to edge-by-edge search heuristics (a
dumbbell must pass through long, low-scoring two-clique-plus-path
intermediate states).
\end{remark}

\begin{remark}[Robustness of the reading]\label{rem:143robust}
Replacing the population variance by the sample variance multiplies the
left-hand side by $k/(k-1) > 1$, so all violations in Table~\ref{tab:143}
persist a fortiori. Reading ``size'' as $n$ instead of $m$ also leaves the
violations intact, since $n < m$ for every instance in the table and the
right-hand side only shrinks. The adversarial audit
(Section~\ref{sec:verification}) additionally reported, at floating-point
level, that replacing average distance by average eccentricity does not
rescue the inequality on these instances. The only readings found to rescue
it are not supported by the source (e.g.\ taking the variance of \emph{all}
eigenvalues).
\end{remark}

The mechanism behind the failure is simple: a dumbbell has exactly two large
positive eigenvalues, close to $t-1$, coming from the cliques, while its
remaining positive eigenvalues stay below $3$; making the path long drives
the variance of the positive spectrum up faster than the ratio
$m/\ell$. This intuition can be made precise, and shows that the conjecture
fails by a factor approaching $2$:

\begin{proposition}\label{prop:143asymp}
For $t \to \infty$ let $p = p(t)$ be integers with $p \to \infty$ and
$p/t \to 0$ (for instance $p = \lfloor t^{1/3} \rfloor$), and let
$G_t = \D(t, p, t)$. Then, with $R = m/\ell$ for either reading
$\ell \in \{\la, \ld\}$,
\[
\frac{\Varp(G_t)}{R(G_t)} \longrightarrow 2,
\qquad\text{and}\qquad
\Varp(G_t) - R(G_t) \;=\; \bigl(1+o(1)\bigr)\,\frac{2t^2}{p+3}
\longrightarrow \infty .
\]
In particular $G_t$ violates Conjecture~143, under both readings, for all
sufficiently large $t$.
\end{proposition}

\begin{proof}
All asymptotic notation refers to $t \to \infty$, and all implied constants
are absolute; recall $n = 2t + p$ and assume throughout $4 \le p \le t$.
Write $A$ for the adjacency matrix of $G_t$ and $A_0$ for that of the
disjoint union $K_t \sqcup P_p \sqcup K_t$, so that $A = A_0 + E$, where $E$
is the adjacency matrix of the two bridge edges. Since $p \ge 2$, the two
bridge edges are vertex-disjoint, so $E$ has eigenvalues $1, 1, -1, -1$ and
$0$ ($n-4$ times); in particular $\lVert E\rVert_2 = 1$.

The spectrum of $A_0$ is classical \cite{BrouwerHaemers2012}: $t-1$ with
multiplicity $2$, $-1$ with multiplicity $2(t-1)$, and
$2\cos\bigl(j\pi/(p+1)\bigr)$ for $j = 1, \dots, p$. Hence for $t \ge 4$ the
two largest eigenvalues of $A_0$ equal $t - 1$, the third largest is
$2\cos\bigl(\pi/(p+1)\bigr) < 2$, and the number of positive eigenvalues of
$A_0$ is $k_0 = 2 + \lfloor p/2 \rfloor$.

By Weyl's perturbation inequality,
$\lvert \lambda_i(A) - \lambda_i(A_0)\rvert \le \lVert E \rVert_2 \le 1$ for
all $i$; therefore
\begin{equation}\label{eq:eigloc}
\lambda_1(A),\, \lambda_2(A) \in [\,t-2,\; t\,],
\qquad
\lambda_i(A) < 3 \quad (i \ge 3).
\end{equation}
Moreover, since $E$ and $-E$ each have exactly two positive eigenvalues, the
Weyl inequalities $\lambda_{i+j-1}(X+Y) \le \lambda_i(X) + \lambda_j(Y)$
(applied with $j = 3$, once to $X = A_0$, $Y = E$ and once to $X = A$,
$Y = -E$) give
$\lambda_{i+2}(A) \le \lambda_i(A_0)$ and
$\lambda_{i+2}(A_0) \le \lambda_i(A)$ for all admissible $i$. Consequently
the number $k$ of (strictly) positive eigenvalues of $A$ satisfies
$k \le k_0 + 2$ --- if $\lambda_i(A) > 0$ for some $i \ge 3$, then
$\lambda_{i-2}(A_0) \ge \lambda_i(A) > 0$, so $i - 2 \le k_0$ --- and
$k \ge k_0 - 2$, since $\lambda_i(A) \ge \lambda_{i+2}(A_0) > 0$ for every
$i \le k_0 - 2$. Hence
\begin{equation}\label{eq:kcount}
\lvert k - k_0 \rvert \le 2, \qquad\text{so}\qquad k = \tfrac{p}{2} + O(1).
\end{equation}
Let $S_1$ and $S_2$ denote the sum and the sum of squares of the positive
eigenvalues of $A$. By \eqref{eq:eigloc} and \eqref{eq:kcount}, the $k - 2$
eigenvalues other than $\lambda_1, \lambda_2$ contribute at most $3(k-2) =
O(p)$ to $S_1$ and at most $9(k-2) = O(p)$ to $S_2$, while
$\lambda_1 + \lambda_2 = 2t + O(1)$ and
$\lambda_1^2 + \lambda_2^2 = 2t^2 + O(t)$. Hence
$S_1 = 2t + O(p)$ and $S_2 = 2t^2 + O(t)$, and
\begin{equation}\label{eq:var}
\Varp(G_t)
= \frac{k S_2 - S_1^2}{k^2}
= \frac{2t^2 (k-2) + O(tp)}{k^2}
= \frac{4t^2}{p}\Bigl(1 + O\bigl(\tfrac1p + \tfrac{p}{t}\bigr)\Bigr),
\end{equation}
using \eqref{eq:kcount} in the last step.

On the other side, $m = t(t-1) + p + 1 = t^2\bigl(1 + O(1/t)\bigr)$. Every
distance in $G_t$ is at most the diameter $p + 3$, and the $(t-1)^2$ pairs
formed by one non-attachment vertex from each clique are at distance exactly
$p+3$; splitting the remaining pairs into within-clique pairs (distance $1$,
at most $t(t-1)$ in total) and pairs involving a path vertex (fewer than
$pn$ of them) gives
\[
(t-1)^2 (p+3) \;\le\; W \;\le\; t^2(p+3) + t(t-1) + pn(p+3),
\]
whence $W = t^2 (p+3)\bigl(1 + O(1/p + p/t)\bigr)$. Since
$n^2 = 4t^2\bigl(1 + O(p/t)\bigr)$ and $n(n-1) = n^2\bigl(1 + O(1/t)\bigr)$,
\begin{equation}\label{eq:rhs}
R(G_t) \;=\; \frac{m\,n^2}{2W} \ \text{ or } \ \frac{m\,n(n-1)}{2W}
\;=\; \frac{2t^2}{p+3}\Bigl(1 + O\bigl(\tfrac1p + \tfrac{p}{t}\bigr)\Bigr).
\end{equation}
Dividing \eqref{eq:var} by \eqref{eq:rhs}:
\[
\frac{\Varp(G_t)}{R(G_t)}
= 2\cdot\frac{p+3}{p}\Bigl(1 + O\bigl(\tfrac1p + \tfrac{p}{t}\bigr)\Bigr)
= 2 + O\bigl(\tfrac1p + \tfrac{p}{t}\bigr) \longrightarrow 2,
\]
since $p \to \infty$ and $p/t \to 0$. Finally
$\Varp - R = R\,(\Varp/R - 1) = (1 + o(1))\, 2t^2/(p+3) \to \infty$, because
$t^2/(p+3) \ge t^2/(t+3) \to \infty$.
\end{proof}

As a numerical illustration (floating point, not part of the certificates):
$\D(40,10,40)$ on $n = 90$ vertices has $m = 1571$, $W = 27625$, $k = 7$,
$\Varp \approx 290.01$ against $m/\la \approx 230.32$ --- a violation margin
of about $59.7$ and a ratio of about $1.26$; for
$\D(800,40,800)$ the ratio reaches about $1.81$.

\section{Conjecture 154: deviation of the eigenvalues}\label{sec:154}

Conjecture~154 reads, verbatim from the decoded WoW source:
\emph{``deviation of eigenvalues $\le$ n / average distance.''} No
surviving source defines ``deviation'' next to the conjecture itself, so the
reading must be fixed from context. WoW uses \emph{variance} (e.g.\
conjectures~38, 143) and \emph{deviation} (e.g.\ conjectures~39, 40, 140,
154, 244, 253) as distinct invariants and spells out ``standard deviation''
for degree sequences (conjectures~27, 812, 813); Favaron, Mah\'eo and
Sacl\'e \cite{FMS1993}, the historical treatment of WoW's
eigenvalue-deviation conjectures, define the eigenvalue deviation (there for
the Laplacian spectrum) by a
root-mean-square formula --- a standard deviation, not a mean absolute
deviation; and the only public machine formalization, the search code of
\cite{RoucairolCazenave2025}, scores the conjecture by the standard
deviation of the adjacency spectrum (its scoring function calls a statistics
helper module that is absent from the published tree; the one deviation
routine the repository does publish, in \texttt{tools/calc.rs}, is a
population standard deviation). We therefore work under the
\textbf{standard-deviation reading}, and claim the refutation \emph{only}
under that reading; Section~\ref{subsec:mad} states exactly what was checked
about the alternative (mean-absolute-deviation) reading.

Under the standard-deviation reading the left-hand side is exactly computable
from $n$ and $m$ alone:

\begin{lemma}\label{lem:dev}
For every graph $G$ with $n$ vertices and $m$ edges, the population standard
deviation of the $n$ adjacency eigenvalues equals
$\dev(G) = \sqrt{2m/n}$.
\end{lemma}

\begin{proof}
$\sum_i \lambda_i = \operatorname{tr} A = 0$ and
$\sum_i \lambda_i^2 = \operatorname{tr} A^2 = \sum_v \deg(v) = 2m$, so the
population variance is $\frac1n \sum_i \lambda_i^2 - \bigl(\frac1n \sum_i
\lambda_i\bigr)^2 = 2m/n$.
\end{proof}

\begin{lemma}\label{lem:integer}
Let $G$ be connected with $n \ge 2$ vertices, $m$ edges and Wiener index
$W$. Under the standard-deviation reading, $G$ violates Conjecture~154
\begin{itemize}\itemsep0pt
\item under the distinct-pairs convention $\ell = \ld$ if and only if
$\;8mW^2 > n^5(n-1)^2$;
\item under the all-entries convention $\ell = \la$ if and only if
$\;8mW^2 > n^7$.
\end{itemize}
\end{lemma}

\begin{proof}
By Lemma~\ref{lem:dev} the conjecture asserts $\sqrt{2m/n} \le n/\ell$. Both
sides are positive, so the inequality is equivalent to its square,
$2m/n \le n^2/\ell^2$. Substituting $\ld = 2W/(n(n-1))$ gives
$2m/n \le n^4(n-1)^2/(4W^2)$, i.e.\ $8mW^2 \le n^5(n-1)^2$; substituting
$\la = 2W/n^2$ gives $8mW^2 \le n^7$. Violation is the strict reverse
inequality. Note $n^7 > n^5(n-1)^2$, so every violation of the $\la$ reading
is automatically one of the $\ld$ reading.
\end{proof}

The verdict is thus a pure integer comparison --- no eigenvalue computation
and no floating point are needed at all. For integers $t \ge 3$, $p \ge 1$,
the \emph{lollipop} $\Lol(t,p)$ is the clique $K_t$ with a pendant path on
$p$ vertices attached by a bridge edge to a clique vertex $v_0$; it is
connected with $n = t + p$ and $m = \binom{t}{2} + p$.

\begin{lemma}\label{lem:wiener}
The Wiener index of the lollipop is
\[
W\bigl(\Lol(t,p)\bigr) = \binom{t}{2} + \sum_{k=1}^{p}\bigl[k + (t-1)(k+1)\bigr]
+ \binom{p+1}{3}
= \binom{t}{2} + \frac{p(p+1)}{2} + (t-1)\,\frac{p(p+3)}{2} + \binom{p+1}{3}.
\]
\end{lemma}

\begin{proof}
Label the path vertices $u_1, \dots, u_p$, with $u_k$ at distance $k$ from
the attachment vertex $v_0$. The $\binom{t}{2}$ clique pairs are at distance
$1$. The pair $\{u_k, v_0\}$ is at distance $k$ and the pairs $\{u_k, w\}$,
for the $t-1$ clique vertices $w \neq v_0$, are at distance $k+1$. The
path-path pairs contribute $\sum_{1 \le i < j \le p} (j - i) =
\binom{p+1}{3}$. Summing gives the first expression; evaluating
$\sum_{k=1}^p k = p(p+1)/2$ and $\sum_{k=1}^p (k+1) = p(p+3)/2$ gives the
second.
\end{proof}

\begin{theorem}\label{thm:154}
Conjecture~154 is false under the standard-deviation reading of
``deviation'', under both conventions for average distance. The integer data
of the three lollipops in Table~\ref{tab:154} are exact;
by Lemma~\ref{lem:integer}, $\Lol(48,70)$ on $n = 118$ vertices violates the
distinct-pairs reading, and $\Lol(50,70)$ on $n = 120$ and $\Lol(72,72)$ on
$n = 144$ vertices violate both readings.
\end{theorem}

\begin{proof}
$m$ and $W$ are evaluated by Lemma~\ref{lem:wiener} (and were re-verified by
breadth-first search over the explicit edge lists); the products and
comparisons in Table~\ref{tab:154} are exact integer arithmetic, reproduced
by two independent implementations (Section~\ref{sec:verification}).
\end{proof}

\begin{table}[ht]
\centering\small
\begin{tabular}{@{}lrrr@{}}
\toprule
 & $\Lol(48,70)$ & $\Lol(50,70)$ & $\Lol(72,72)$ \\
\midrule
$n$ & $118$ & $120$ & $144$ \\
$m$ & $1198$ & $1295$ & $2628$ \\
$W$ & $180\,853$ & $186\,060$ & $259\,080$ \\
$8mW^2$        & $313\,471\,628\,124\,656$ & $358\,645\,832\,496\,000$ & $1\,411\,182\,313\,113\,600$ \\
$n^5(n-1)^2$   & $313\,171\,159\,328\,352$ & $352\,370\,995\,200\,000$ & $1\,266\,148\,181\,016\,576$ \\
$n^7$          & $318\,547\,390\,056\,832$ & $358\,318\,080\,000\,000$ & $1\,283\,918\,464\,548\,864$ \\
$8mW^2 - n^5(n-1)^2$ & $+300\,468\,796\,304$ & $+6\,274\,837\,296\,000$ & $+145\,034\,132\,097\,024$ \\
$8mW^2 - n^7$        & $-5\,075\,761\,932\,176$ & $+327\,752\,496\,000$ & $+127\,263\,848\,564\,736$ \\
violates       & $\ld$ only & both & both \\
\bottomrule
\end{tabular}
\caption{Integer certificates for Theorem~\ref{thm:154}
(Lemma~\ref{lem:integer}: violation of the $\ld$ reading iff
$8mW^2 > n^5(n-1)^2$, of the $\la$ reading iff $8mW^2 > n^7$).}
\label{tab:154}
\end{table}

For orientation: $\dev(\Lol(72,72)) = \sqrt{73/2} = 6.041522\ldots$ against
$n/\la = n^3/(2W) = 62208/10795 = 5.762667\ldots$ and
$n/\ld = n^2(n-1)/(2W) = 61776/10795 = 5.722649\ldots$; in the squared
comparison of Lemma~\ref{lem:integer} the threshold is exceeded by about
$11\%$ ($\ld$) and $10\%$ ($\la$). The whole refutation can be re-verified
in well under a minute; the following self-contained check (quoted verbatim
from the certificate bundle) needs nothing beyond a Python interpreter:

\begin{quote}
\begin{verbatim}
t, p = 72, 72
n, m = t+p, t*(t-1)//2 + p
W = t*(t-1)//2 + sum(k + (t-1)*(k+1) for k in range(1, p+1)) \
    + (p+1)*p*(p-1)//6
assert 8*m*W*W > n**5 * (n-1)**2     # distinct-pairs convention
assert 8*m*W*W > n**7                # all-n^2-entries convention
print("Graffiti 154 (std-dev reading) refuted:", n, m, W)
\end{verbatim}
\end{quote}

The failure is again structural: the standard deviation $\sqrt{2m/n}$ is
unbounded along lollipops with $p = t$, while $n/\ell$ stays bounded.

\begin{proposition}\label{prop:154asymp}
For every integer $t \ge 71$, the lollipop $\Lol(t,t)$ on $n = 2t$ vertices
violates Conjecture~154 (standard-deviation reading) under both conventions
for average distance, and
\[
\frac{\dev\bigl(\Lol(t,t)\bigr)}{\,n/\ell\,} \;\ge\;
\frac{\sqrt{(t+1)/2}}{6} \;\longrightarrow\; \infty
\qquad (\ell \in \{\la, \ld\}).
\]
\end{proposition}

\begin{proof}
Here $m = \binom{t}{2} + t = t(t+1)/2$, so $\dev = \sqrt{2m/n} =
\sqrt{(t+1)/2}$ exactly, by Lemma~\ref{lem:dev}. Specializing
Lemma~\ref{lem:wiener} to $p = t$:
\begin{align*}
W &= \frac{t(t-1)}{2} + \frac{t(t+1)}{2} + \frac{t(t-1)(t+3)}{2}
+ \frac{(t+1)t(t-1)}{6}\\
&= \frac{t\,\bigl(6t + 3(t-1)(t+3) + (t^2-1)\bigr)}{6}
= \frac{t\,(2t^2 + 6t - 5)}{3}.
\end{align*}
Hence
\[
\frac{n}{\ld} = \frac{n^2(n-1)}{2W} = \frac{6t(2t-1)}{2t^2+6t-5}
\qquad\text{and}\qquad
\frac{n}{\la} = \frac{n^3}{2W} = \frac{12t^2}{2t^2+6t-5}.
\]
Both are smaller than $6$ for every $t \ge 1$: indeed
$6t(2t-1) < 6(2t^2+6t-5)$ amounts to $30 < 42t$, and
$12t^2 < 6(2t^2+6t-5)$ amounts to $30 < 36t$. For $t \ge 71$ we have
$\dev = \sqrt{(t+1)/2} \ge 6 > n/\ell$ for both conventions, which is the
asserted (strict) violation, and the ratio bound follows.
\end{proof}

\begin{remark}[Minimality within the family]\label{rem:154min}
An exhaustive integer-arithmetic scan of all lollipops $\Lol(t,p)$ with
$t \ge 3$, $p \ge 1$ and $n = t + p \le 130$ (re-run, like all computations
quoted here, while preparing this note; the independent audit re-ran it
allowing the degenerate values $t = 1, 2$ as well) shows: no lollipop with
$n \le 117$ violates either reading; at $n = 118$ there are exactly four
violators of the $\ld$ reading, $\Lol(48,70)$, $\Lol(49,69)$, $\Lol(50,68)$,
$\Lol(51,67)$, none violating the $\la$ reading; at $n = 119$ exactly seven
lollipops, $\Lol(t, 119-t)$ for $47 \le t \le 53$, violate the $\ld$ reading
and again none the $\la$ reading; and at $n = 120$ exactly two
lollipops, $\Lol(50,70)$ and $\Lol(51,69)$, violate both readings. Since
this scan is exact, these statements are certificates, not numerics ---
but they concern the lollipop family only. The minimum order of a
counterexample over all connected graphs is not determined here: a smaller,
non-lollipop violator of any order up to $117$ is not excluded --- the prior
search for this conjecture \cite{RoucairolCazenave2025} reached size~50, but
it was a heuristic search, not an exhaustive enumeration.
\end{remark}

\subsection{The mean-absolute-deviation reading}\label{subsec:mad}

If Graffiti's ``deviation'' meant the \emph{mean absolute deviation} (MAD)
rather than the standard deviation, the instances above would prove nothing:
we record here, for completeness, exactly what was checked. Since the adjacency
spectrum has mean zero, the MAD equals
$\frac1n \sum_i \lvert\lambda_i\rvert = \mathcal{E}(G)/n$, where
$\mathcal{E}(G)$ is the graph energy. For each
$G \in \{\Lol(48,70), \Lol(50,70), \Lol(72,72)\}$ it was \emph{proved in
exact arithmetic} that
\[
\frac{\mathcal{E}(G)}{n} \;<\; \frac{n}{\ell}
\qquad \text{for both } \ell \in \{\la, \ld\},
\]
i.e.\ these lollipops do \emph{not} violate the MAD reading. Two independent
exact routes were used: (i) a Sturm-certified rational enclosure of
$\mathcal{E}(G)$ of width below $3 \times 10^{-17}$, computed from the
integer characteristic polynomial of the $(p+2)\times(p+2)$ quotient matrix
of the equitable partition $\{v_0\}, V(K_t)\setminus\{v_0\}, \{u_1\}, \dots,
\{u_p\}$, together with the eigenvalue $-1$ of multiplicity $t-2$ (both
verified by explicit eigenvectors and the two trace identities); and (ii) an
exact-rational application of the Koolen--Moulton
energy bound \cite{KoolenMoulton2001}. Numerically, the certified energies
are $\mathcal{E} = 183.02566963\ldots$, $187.02616866\ldots$,
$233.57628649\ldots$ respectively, giving MAD values
$1.5510\ldots$, $1.5585\ldots$, $1.6220\ldots$ against right-hand sides
between $4.5039\ldots$ and $5.7626\ldots$ --- sign gaps of $+2.95$ to
$+4.14$, certified. The refutation of Conjecture~154 in this note is
therefore claimed \emph{only} under the standard-deviation reading, which is
the reading of \cite{FMS1993} and of the only known machine formalization
\cite{RoucairolCazenave2025}; the residual interpretive risk is thereby
documented, quantified, and confined. (For the sample standard deviation,
which exceeds the population standard deviation by the factor
$\sqrt{n/(n-1)} > 1$, the violations hold a fortiori.)

\section{Verification}\label{sec:verification}

Both refutations are packaged as machine-checkable certificate bundles, with
no floating point on any accept path. Each result is certified by two
algorithmically independent exact-arithmetic routes: for Conjecture~143 the
two routes live inside a single clean-room checker, while Conjecture~154 has
two independently written checkers.

\paragraph{Statement provenance.}
The conjecture statements were frozen from the compilation of
\emph{Written on the Wall} archived as \texttt{wow-july2004.pdf} \cite{WOW}
--- the same file used by \cite{RoucairolCazenave2025}, pinned here by the
checksum below. That PDF uses Type-3 bitmap fonts without
ToUnicode maps; the statements were recovered by a glyph-decoding script
validated against the document title and the surrounding conjectures; the
statement of Conjecture~154 was additionally confirmed by pixel-level
rendering of the relevant pages during the adversarial audit.

\paragraph{Conjecture 143.}
The certificate lists, for each of the five dumbbells of
Table~\ref{tab:143}, the explicit edge list, the integers $n, m, W, k$, the
exact rational right-hand sides, a rational enclosure of $\Varp$ of width
below $10^{-19}$, and exact rational lower bounds for every claimed margin.
The checker (written clean-room, independently of the discovery pipeline)
verifies every instance by two algorithmically independent exact routes and
accepts only if both certify each claimed strict inequality and, for each
reading not claimed (the dashes of Table~\ref{tab:143}), that the certified
enclosure of $\Varp$ lies strictly below the right-hand side: route A
computes the integer
characteristic polynomial symbolically and isolates its real roots exactly
(isolation width $2^{-70}$); route B is restricted to the Python standard
library (exact fractions only) and recomputes the characteristic polynomial
by the Berkowitz algorithm, checks it against an equitable-partition quotient
via the exact integer identity
$\operatorname{char}_A(x) = \operatorname{char}_B(x)\,(x+1)^{t_1+t_2-4}$,
isolates roots by Sturm chains with bisection refinement to width $2^{-70}$,
and evaluates the variance in rational interval arithmetic. Cross-checks
include identical characteristic polynomials, trace and edge-count
coefficient identities, agreement of $k$, intersection of the two variance
enclosures, connectivity, and recomputation of $W$, $m$ and the right-hand
sides from the edge lists. Ten targeted certificate corruptions (mutation
tests) are all rejected, the pristine certificate accepted. An independent
audit recomputed all five instances with separate numerical software and
confirmed every certificate field, and an orchestrator-level script rebuilt
$\D(20,8,20)$ from the construction description alone (rather than from the
certified edge list) and reproduced both violations at floating-point
precision.

\paragraph{Conjecture 154.}
Here the accept path is integer-only (Lemma~\ref{lem:integer}). Checker A
rebuilds the three lollipops from scratch, recomputes $W$ both by
breadth-first search and by the closed form of Lemma~\ref{lem:wiener}
(demanding agreement), evaluates the integer violation forms, re-runs the
exhaustive $n \le 130$ family scan of Remark~\ref{rem:154min}, and proves the
MAD non-violation of Section~\ref{subsec:mad} in exact arithmetic by the two
routes described there; the spectral side-checks (trace identities at double
and 60-digit precision, agreement of the certified spectrum with a numerical
eigensolver) are advisory only. Checker B was written independently by an
adversarial referee agent from a different model family (Codex, GPT-5.5)
and reproduces the
integer inequalities, the margins, and the $n \le 117$ non-violation from
its own reconstruction. Six mutation tests (e.g.\ claiming a violation at
$\Lol(47,70)$, corrupting $W$ by one, claiming the $\la$ reading at
$\Lol(48,70)$) are all rejected. An orchestrator-level script again rebuilds
$\Lol(72,72)$ from the construction description and re-asserts both integer
inequalities.

\paragraph{Independent recomputation for this note.}
In June 2026, while preparing this note, both checkers and the mutation
tests were re-run (verdicts: \texttt{ACCEPT} and \texttt{ALL CHECKS PASS},
with all mutation tests rejected); these re-runs certify every entry of
Tables~\ref{tab:143} and~\ref{tab:154}, the lollipop family scan of
Remark~\ref{rem:154min}, and the mean-absolute-deviation bounds of
Section~\ref{subsec:mad}. The dumbbell family scan of
Remark~\ref{rem:143min} was re-executed, and the headline quantities --- the
values of Table~\ref{tab:143} and the numerical illustrations
$\D(40,10,40)$, $\D(800,40,800)$ and $\Lol(t,t)$ --- were additionally
recomputed by separate scripts written for the purpose.

\paragraph{Data availability.}
The two certificate bundles --- machine-readable certificates (JSON),
all checkers, the mutation-test suites, the frozen source PDF of
\emph{Written on the Wall} with its decoding script, the archived search code
of \cite{RoucairolCazenave2025} used to fix the formalization conventions,
and full verification logs --- are available at
\url{https://doi.org/10.5281/zenodo.20673865} and
\url{https://github.com/demonstrandum-research/artifacts}. The frozen source PDF has SHA-256 checksum
\begin{center}
\texttt{\small d2c779d2c28418b30ab1b4f84bf7112c0e000bca1f2d6a3c48d0baba78af7733}.
\end{center}

\begin{thebibliography}{10}

\bibitem{AouchicheHansen2010}
M.~Aouchiche and P.~Hansen,
\emph{A survey of automated conjectures in spectral graph theory},
Linear Algebra Appl.\ \textbf{432} (2010), 2293--2322.

\bibitem{BDF1995}
T.~L. Brewster, M.~J. Dinneen, and V.~Faber,
\emph{A computational attack on the conjectures of Graffiti: new
counterexamples and proofs},
Discrete Math.\ \textbf{147} (1995), 35--55.

\bibitem{BrouwerHaemers2012}
A.~E. Brouwer and W.~H. Haemers,
\emph{Spectra of Graphs},
Universitext, Springer, New York, 2012.

\bibitem{DeLaVina2005}
E.~DeLaVi\~na,
\emph{Some history of the development of Graffiti},
in: Graphs and Discovery, DIMACS Ser.\ Discrete Math.\ Theoret.\ Comput.\
Sci.\ \textbf{69}, Amer.\ Math.\ Soc., Providence, RI, 2005, pp.~81--118.

\bibitem{Fajtlowicz1988}
S.~Fajtlowicz,
\emph{On conjectures of Graffiti},
Discrete Math.\ \textbf{72} (1988), 113--118.

\bibitem{WOW}
S.~Fajtlowicz,
\emph{Written on the Wall (conjectures of Graffiti)},
working document, University of Houston, regularly updated; compilation
archived by Roucairol and Cazenave as \texttt{wow-july2004.pdf} (the PDF
metadata records a creation date of 26 July 2006). Archived copy (accessed
2026-06-11):
\url{https://raw.githubusercontent.com/RoucairolMilo/refutation-COCOON2022/master/wow-july2004.pdf}.

\bibitem{FMS1993}
O.~Favaron, M.~Mah\'eo, and J.-F. Sacl\'e,
\emph{Some eigenvalue properties in graphs (conjectures of Graffiti---II)},
Discrete Math.\ \textbf{111} (1993), 197--220.

\bibitem{KoolenMoulton2001}
J.~H. Koolen and V.~Moulton,
\emph{Maximal energy graphs},
Adv.\ in Appl.\ Math.\ \textbf{26} (2001), 47--52.

\bibitem{RoucairolCazenave2022}
M.~Roucairol and T.~Cazenave,
\emph{Refutation of spectral graph theory conjectures with Monte Carlo
search},
in: Computing and Combinatorics (COCOON 2022), Lecture Notes in Comput.\
Sci.\ \textbf{13595}, Springer, 2022, pp.~162--176.

\bibitem{RoucairolCazenave2025}
M.~Roucairol and T.~Cazenave,
\emph{Refutation of spectral graph theory conjectures with search
algorithms},
presented at the Workshop on Intelligent Agents in Science and Engineering
(IASE-2025) at ECAI 2025, Bologna, October 2025; arXiv:2409.18626.

\end{thebibliography}

\end{document}
